\section{Abstract}
Der Einsatz von Large Language Models (LLMs) in industriellen Automatisierungsprozessen – insbesondere an der Schnittstelle zu ERP-Systemen – ermöglicht eine stark vereinfachte Übersetzung natürlichsprachlicher Anforderungen in ausführbare Workflows. Gleichzeitig entstehen jedoch neue Risiken: nicht-deterministisches Verhalten, Halluzinationen, fehlende Nachvollziehbarkeit sowie unzureichende Validierung der erzeugten Ergebnisse. Dadurch bleiben gerade jene Anwendungsfälle ausgeschlossen, in denen Prozesssicherheit, Auditierbarkeit und korrekte ERP-Transaktionen zwingend sind (z.\,B. Stammdatenänderungen, Bestellfreigaben, Buchungen).

Diese Masterarbeit untersucht, wie Quality-Gate-Mechanismen so in eine LLM-Pipeline integriert werden können, dass Korrektheit und Verlässlichkeit KI-generierter Automatisierungsartefakte \textbf{nicht erst am Ende}, sondern \textbf{kontinuierlich im laufenden Erstellungsprozess} abgesichert werden. Im Fokus steht ein praxisnahes Szenario, in dem Kundenanforderungen schrittweise von unstrukturierter Sprache in eine strukturierte, maschinenlesbare Automatisierung überführt und mit ERP-Funktionen verknüpft werden. Ausgangspunkt ist ein unpräziser Kundenwunsch, der durch ein geführtes Eingabeformat und formale Constraints präzisiert wird. Darauf aufbauend generiert ein LLM-gestütztes System iterativ einen vollständigen Automatisierungsworkflow als JSON-Repräsentation für die Plattform n8n.

Die Quality Gates werden dabei als explizite, versionierte Prüfschritte modelliert und automatisch erzeugt bzw.\ abgeleitet: (1) formale Validierung (Schema/Typen/Constraints), (2) semantische Konsistenzprüfungen und Traceability (Anforderung $\rightarrow$ Schritt $\rightarrow$ Parameter), (3) Schnittstellen- und Vertragsprüfungen zu ERP-APIs (z.\,B. erlaubte Felder, Idempotenz, Berechtigungen), (4) simulations- bzw.\ testbasierte Verifikation auf Beispiel- und Grenzfälle sowie (5) kontrollierte Freigabe mit revisionssicherem Protokoll. Ziel ist eine Pipeline-Architektur, die LLM-Generierung und Quality Gates so koppelt, dass fehlerhafte Teilartefakte früh erkannt, repariert und erst nach bestandenen Gates in den nächsten Schritt übernommen werden.